\section{Time and space complexity}\label{sec:complexity}

We analyze Algorithm~\ref{alg:astar} with a consistent heuristic, where $|V|$
and $|E|$ are the numbers of vertices and edges. We assume that $Q$ is a binary
heap, in which each queue operation takes $O(\log |V|)$ time, and that $h(v)$ is
computed in constant time, as for the distance heuristics.

\subsection{Time}\label{sec:time}

By Theorem~\ref{thm:consistent}, every vertex is extracted at most once, so
every edge is examined at most twice (Table~\ref{tab:time}).

\begin{table}[htbp]
  \centering
  \caption{Running time of Algorithm~\ref{alg:astar} with a binary
    heap.}\label{tab:time}
  \renewcommand{\arraystretch}{1.25}%
  \begin{tabular}{@{}P{4.4cm}P{4.6cm}P{2.2cm}P{2.6cm}@{}}
    \toprule
    Operation & Number of executions & Cost of each & Total\\
    \midrule
    initialization of $g$ and $\pi$ & once per vertex & $O(1)$ & $O(|V|)$\\
    \textsc{Extract-Min} & at most once per vertex & $O(\log|V|)$
      & $O(|V|\log|V|)$\\
    examination of an edge & at most twice per edge & $O(1)$ & $O(|E|)$\\
    \textsc{Insert-or-Decrease-Key} & at most once per examination of an edge
      & $O(\log|V|)$ & $O(|E|\log|V|)$\\
    \bottomrule
  \end{tabular}
\end{table}

Summing the totals gives the running time
\begin{equation}\label{eq:time}
  O(|V|) + O(|V| \log |V|) + O(|E|) + O(|E| \log |V|)
  = O\big((|V| + |E|) \log |V|\big).
\end{equation}

This is also the bound for Dijkstra's algorithm, and a Fibonacci heap reduces
both to $O(|E| + |V| \log |V|)$~\cite{cormen2022clrs}. In the worst case A* is
no faster than Dijkstra's algorithm, since the zero heuristic is consistent; its
advantage lies in the number of vertices expanded
(Section~\ref{sec:efficiency}).

With an inconsistent heuristic, A* with reopening can expand a vertex many
times; Martelli~\cite{martelli1977} gave graphs on which the number of
expansions grows exponentially.

\subsection{Space}\label{sec:space}

Each item of Table~\ref{tab:space} holds at most one entry per vertex, so the
memory is $O(|V|)$.

\begin{table}[htbp]
  \centering
  \caption{Memory used by Algorithm~\ref{alg:astar}.}\label{tab:space}
  \renewcommand{\arraystretch}{1.25}%
  \begin{tabular}{@{}P{5cm}P{6cm}P{2.4cm}@{}}
    \toprule
    Data & Size & Total\\
    \midrule
    values $g$ and predecessors $\pi$ & one entry per vertex & $O(|V|)$\\
    closed set $S$ & at most every vertex & $O(|V|)$\\
    priority queue $Q$ & at most one entry per vertex, with
      \textsc{Decrease-Key} & $O(|V|)$\\
    reconstructed path & at most $|V|$ vertices & $O(|V|)$\\
    \midrule
    \textbf{Total} & & $O(|V|)$\\
    \bottomrule
  \end{tabular}
\end{table}

Implementations that insert a second queue entry instead of decreasing a key
need $O(|V| + |E|)$ space. On large graphs memory, not time, is usually the
limit, because A* stores every vertex it reaches~\cite{russell2020aima}.

\subsection{Comparison of the three algorithms}\label{sec:complexity-compare}

In greedy best-first search every vertex enters $Q$ once and no key decreases,
which gives the smallest bound in Table~\ref{tab:complexity}. The worst-case
bounds do not separate A* from Dijkstra's algorithm.

\begin{table}[htbp]
  \centering
  \caption{Worst-case complexity of the three algorithms with a binary heap and
    a consistent heuristic.}\label{tab:complexity}
  \renewcommand{\arraystretch}{1.25}%
  \begin{tabular}{@{}P{2.6cm}P{3.6cm}P{3.6cm}P{3.6cm}@{}}
    \toprule
     & \textcolor{algdijkstra}{\textbf{Dijkstra's algorithm}}
     & \textcolor{alggreedy}{\textbf{Greedy best-first search}}
     & \textcolor{algastar}{\textbf{A* search}}\\
    \midrule
    Time  & $O((|V|{+}|E|)\log|V|)$ & $O(|E| + |V|\log|V|)$
      & $O((|V|{+}|E|)\log|V|)$\\[2pt]
    Space & $O(|V|)$ & $O(|V|)$ & $O(|V|)$\\[2pt]
    Shortest path & \textcolor{algastar}{guaranteed}
      & \textcolor{cbRed}{not guaranteed} & \textcolor{algastar}{guaranteed}\\
    \bottomrule
  \end{tabular}
\end{table}
